% ============================================================
%  MTH636 -- Game Theory | Lecture 12
%  Equivalence of Minimax and Nash, Indifference Principle
% ============================================================
\documentclass[12pt, a4paper]{article}

% ============================================================
%  MTH636 -- Game Theory
%  Notes Preamble — Color-Coded, TikZ-Enhanced
% ============================================================

% --- Core packages ---
\usepackage[T1]{fontenc}
\usepackage[utf8]{inputenc}
\usepackage{lmodern}
\usepackage[margin=2.5cm, top=2.8cm, bottom=2.8cm]{geometry}
\usepackage{microtype}

% --- Math ---
\usepackage{amsmath, amssymb, amsthm, mathtools}
\usepackage{bm}        % bold math
\usepackage{cancel}    % strike-through in math

% --- Color & Boxes ---
\usepackage[dvipsnames, svgnames, x11names, table]{xcolor}
\usepackage[most]{tcolorbox}
\tcbuselibrary{breakable, skins, theorems}

% --- TikZ & Diagrams ---
\usepackage{tikz}
\usepackage{tikz-qtree}
\usetikzlibrary{arrows.meta, positioning, shapes.geometric, fit, backgrounds,
                decorations.pathreplacing, calc, matrix, trees}
\usepackage{pgfplots}
\pgfplotsset{compat=1.18}

% --- Tables ---
\usepackage{booktabs}
\usepackage{multirow}
\usepackage{array}
\usepackage{tabularx}
\usepackage{makecell}

% --- Typography & Layout ---
\usepackage{enumitem}
\usepackage{parskip}
\usepackage{fancyhdr}
\usepackage{titlesec}
\usepackage{titletoc}
\usepackage{hyperref}
\usepackage{graphicx}
\usepackage{float}
\usepackage{caption}

% --- Color Palette ---
\definecolor{defblue}{HTML}{1A5276}      % Definition header
\definecolor{defbluebg}{HTML}{D6EAF8}    % Definition background
\definecolor{thmgreen}{HTML}{1E8449}     % Theorem header
\definecolor{thmbg}{HTML}{D5F5E3}        % Theorem background
\definecolor{exorange}{HTML}{B7770D}     % Example header
\definecolor{exbg}{HTML}{FEF9E7}         % Example background
\definecolor{warnred}{HTML}{922B21}      % Warning/Key Insight header
\definecolor{warnbg}{HTML}{FADBD8}       % Warning background
\definecolor{noticegray}{HTML}{2C3E50}   % Note header
\definecolor{noticebg}{HTML}{EBF5FB}     % Note background
\definecolor{taketawayyellow}{HTML}{7D6608} % Key Takeaway header
\definecolor{takeawaybg}{HTML}{FDFEFE}   % Takeaway background
\definecolor{accentpurple}{HTML}{6C3483} % Accent / proof
\definecolor{lightgray}{HTML}{F2F3F4}    % shading

% ============================================================
%  TCOLORBOX ENVIRONMENTS
% ============================================================

% -- Definition (Blue) --
\newtcolorbox{defbox}[2][]{
  enhanced, breakable,
  colback=defbluebg, colframe=defblue, coltitle=white,
  fonttitle=\bfseries\sffamily,
  title={Definition: #2},
  attach boxed title to top left={yshift=-2mm, xshift=4mm},
  boxed title style={colback=defblue, rounded corners},
  arc=4pt, boxrule=1.2pt,
  left=6pt, right=6pt, top=8pt, bottom=6pt,
  #1
}

% -- Theorem / Proposition (Green) --
\newtcolorbox{thmbox}[2][]{
  enhanced, breakable,
  colback=thmbg, colframe=thmgreen, coltitle=white,
  fonttitle=\bfseries\sffamily,
  title={Theorem: #2},
  attach boxed title to top left={yshift=-2mm, xshift=4mm},
  boxed title style={colback=thmgreen, rounded corners},
  arc=4pt, boxrule=1.2pt,
  left=6pt, right=6pt, top=8pt, bottom=6pt,
  #1
}

% -- Proposition (Green variant) --
\newtcolorbox{propbox}[2][]{
  enhanced, breakable,
  colback=thmbg, colframe=thmgreen!70!black, coltitle=white,
  fonttitle=\bfseries\sffamily,
  title={Proposition: #2},
  attach boxed title to top left={yshift=-2mm, xshift=4mm},
  boxed title style={colback=thmgreen!70!black, rounded corners},
  arc=4pt, boxrule=1.2pt,
  left=6pt, right=6pt, top=8pt, bottom=6pt,
  #1
}

% -- Example (Orange/Yellow) --
\newtcolorbox{exbox}[2][]{
  enhanced, breakable,
  colback=exbg, colframe=exorange, coltitle=white,
  fonttitle=\bfseries\sffamily,
  title={Example: #2},
  attach boxed title to top left={yshift=-2mm, xshift=4mm},
  boxed title style={colback=exorange, rounded corners},
  arc=4pt, boxrule=1.2pt,
  left=6pt, right=6pt, top=8pt, bottom=6pt,
  #1
}

% -- Key Insight / Warning (Red) --
\newtcolorbox{warnbox}[2][]{
  enhanced, breakable,
  colback=warnbg, colframe=warnred, coltitle=white,
  fonttitle=\bfseries\sffamily,
  title={Key Insight: #2},
  attach boxed title to top left={yshift=-2mm, xshift=4mm},
  boxed title style={colback=warnred, rounded corners},
  arc=4pt, boxrule=1.2pt,
  left=6pt, right=6pt, top=8pt, bottom=6pt,
  #1
}

% -- Note (Gray-Blue) --
\newtcolorbox{notebox}[1][]{
  enhanced, breakable,
  colback=noticebg, colframe=noticegray!60, coltitle=white,
  fonttitle=\bfseries\sffamily,
  title={Note},
  attach boxed title to top left={yshift=-2mm, xshift=4mm},
  boxed title style={colback=noticegray!60, rounded corners},
  arc=4pt, boxrule=0.8pt,
  left=6pt, right=6pt, top=8pt, bottom=6pt,
  #1
}

% -- Key Takeaway --
\newtcolorbox{takebox}[1][]{
  enhanced, breakable,
  colback=takeawaybg, colframe=taketawayyellow, coltitle=white,
  fonttitle=\bfseries\sffamily,
  title={Key Takeaway},
  attach boxed title to top left={yshift=-2mm, xshift=4mm},
  boxed title style={colback=taketawayyellow, rounded corners},
  arc=4pt, boxrule=1.2pt,
  left=6pt, right=6pt, top=8pt, bottom=6pt,
  #1
}

% -- Intuition --
\newtcolorbox{intubox}[1][]{
  enhanced, breakable,
  colback=accentpurple!8, colframe=accentpurple!60, coltitle=white,
  fonttitle=\bfseries\sffamily,
  title={Intuition},
  attach boxed title to top left={yshift=-2mm, xshift=4mm},
  boxed title style={colback=accentpurple!80, rounded corners},
  arc=4pt, boxrule=0.8pt,
  left=6pt, right=6pt, top=8pt, bottom=6pt,
  #1
}

% ============================================================
%  PAYOFF MATRIX HELPER
%  Usage: inside tabular, use \payoff{a}{b} for (a,b)
% ============================================================
\newcommand{\payoff}[2]{$(#1,\;#2)$}
\newcommand{\nash}[2]{\underline{$#1$},\underline{$#2$}}  % underline NE payoffs
\newcommand{\Nset}{\mathbb{N}}
\newcommand{\Rset}{\mathbb{R}}
\newcommand{\Sset}{\mathcal{S}}
\newcommand{\Ui}{u_i}

% ============================================================
%  SECTION FORMATTING
% ============================================================
\titleformat{\section}
  {\large\bfseries\sffamily\color{defblue}}
  {\thesection}{1em}{}[\titlerule]

\titleformat{\subsection}
  {\normalsize\bfseries\sffamily\color{thmgreen!80!black}}
  {\thesubsection}{1em}{}

\titleformat{\subsubsection}
  {\normalsize\itshape\sffamily\color{exorange!80!black}}
  {\thesubsubsection}{1em}{}

% ============================================================
%  HEADER / FOOTER
% ============================================================
\pagestyle{fancy}
\fancyhf{}
\fancyhead[L]{\sffamily\small\color{defblue} MTH636 -- Game Theory}
\fancyhead[R]{\sffamily\small\color{defblue} \leftmark}
\fancyfoot[C]{\sffamily\small\thepage}
\renewcommand{\headrulewidth}{0.4pt}

% ============================================================
%  HYPERREF SETUP
% ============================================================
\hypersetup{
  colorlinks=true,
  linkcolor=defblue,
  urlcolor=thmgreen,
  citecolor=accentpurple,
  bookmarksnumbered=true,
  pdftitle={MTH636 Game Theory Notes},
  pdfauthor={MTH636 Student Notes}
}

% ============================================================
%  MISC HELPERS
% ============================================================
\newcommand{\BR}{\text{BR}}        % Best Response
\newcommand{\NE}{\text{NE}}        % Nash Equilibrium
\newcommand{\SPNE}{\text{SPNE}}    % Subgame Perfect NE
\newcommand{\Dom}{\text{dom}}      % Dominance
\newcommand{\mixs}[1]{\sigma_{#1}} % Mixed strategy
\newcommand{\pures}[1]{s_{#1}}     % Pure strategy


% ============================================================
\begin{document}

% Title
\begin{center}
  {\LARGE\bfseries\sffamily\color{defblue} MTH636 -- Game Theory}\\[6pt]
  {\Large\sffamily Lecture 12: NE and the Indifference Principle}\\[4pt]
\end{center}
\vspace{0.5cm}
\hrule height 1.5pt
\vspace{0.3cm}

% ============================================================
\section{Minimax Strategies and Nash Equilibrium}

For general two-player zero-sum games, finding max-min strategies can be framed as a Linear Programming problem (see Solan's book, Section 5.2.5). There is a profound connection between max-min strategies and Nash Equilibrium.

\begin{thmbox}{Equivalence Theorem for Zero-Sum Games}
  For a two-player zero-sum game, a strategy profile $(\sigma_1^*, \sigma_2^*)$ is a max-min strategy profile if and only if $(\sigma_1^*, \sigma_2^*)$ is a Nash Equilibrium.
\end{thmbox}

\textbf{Proof Idea:}
If $V$ is the value of the game, a max-min profile means $U(\sigma_1^*, \sigma_2) \geq V$ for all $\sigma_2$ and $U(\sigma_1, \sigma_2^*) \leq V$ for all $\sigma_1$. 
Setting $\sigma_2 = \sigma_2^*$ and $\sigma_1 = \sigma_1^*$ yields $U(\sigma_1^*, \sigma_2^*) = V$.
Therefore:
\[
  U(\sigma_1^*, \sigma_2^*) = V \geq U(\sigma_1, \sigma_2^*) \quad \forall \sigma_1 \in \Sigma_1
\]
\[
  U(\sigma_1^*, \sigma_2^*) = V \leq U(\sigma_1^*, \sigma_2) \quad \forall \sigma_2 \in \Sigma_2
\]
This precisely satisfies the definition of a Nash Equilibrium. The proof in the other direction follows a similar bounding logic to show that the existence of an NE implies $\underline{V} = \overline{V} = V$.

\begin{thmbox}{Security Level Theorem}
  Suppose $\sigma^*$ is a Nash Equilibrium of any strategic form game (not necessarily zero-sum). Then every player achieves at least their max-min value (security level):
  \[
    U_i(\sigma^*) \geq \underline{V}_i \quad \forall i \in N
  \]
\end{thmbox}

% ============================================================
\section{Computing NE via Best Response}

Let's compute the mixed Nash Equilibrium for a non-zero-sum game using best response correspondences.

\begin{exbox}{Battle of the Sexes}
  \begin{center}
  \renewcommand{\arraystretch}{1.3}
  \begin{tabular}{c|cc}
    $P_1 \setminus P_2$ & L ($y$) & R ($1-y$) \\
    \hline
    T ($x$) & (2,1) & (0,0) \\
    B ($1-x$) & (0,0) & (1,2) \\
  \end{tabular}
  \end{center}
\end{exbox}

\textbf{Player 1's Expected Utility:}
\begin{align*}
  U_1(x, y) &= x[2y + 0] + (1-x)[0 + 1(1-y)] \\
  &= 2xy + 1 - y - x + xy = x(3y - 1) + (1-y)
\end{align*}
To maximize this, Player 1 looks at the coefficient of $x$, which is $(3y - 1)$.
\begin{itemize}
    \item If $3y - 1 > 0 \implies y > 1/3$, P1 wants $x$ to be as large as possible ($x=1$).
    \item If $3y - 1 < 0 \implies y < 1/3$, P1 wants $x$ to be as small as possible ($x=0$).
    \item If $3y - 1 = 0 \implies y = 1/3$, P1 is indifferent, any $x \in [0,1]$ is a best response.
\end{itemize}

\textbf{Player 2's Expected Utility:}
\begin{align*}
  U_2(x, y) &= y[1x + 0] + (1-y)[0 + 2(1-x)] \\
  &= xy + 2 - 2x - 2y + 2xy = y(3x - 2) + 2 - 2x
\end{align*}
Looking at the coefficient of $y$, which is $(3x - 2)$:
\begin{itemize}
    \item If $x > 2/3$, P2 wants $y=1$.
    \item If $x < 2/3$, P2 wants $y=0$.
    \item If $x = 2/3$, P2 is indifferent, any $y \in [0,1]$ is a best response.
\end{itemize}

If we graph these two best response correspondences, they intersect at three points:
1. $(x=0, y=0)$ $\implies$ Pure NE: $(B, R)$
2. $(x=1, y=1)$ $\implies$ Pure NE: $(T, L)$
3. $(x=2/3, y=1/3)$ $\implies$ \textbf{Mixed NE:} P1 plays T with prob $2/3$, P2 plays L with prob $1/3$.

% ============================================================
\section{The Indifference Principle}

Notice in the previous example that in the mixed equilibrium, each player makes the \textit{other} player indifferent between their pure strategies. This is a fundamental property of all mixed strategy Nash Equilibria.

\begin{thmbox}{Support of a Mixed Strategy}
  Suppose $\sigma^*$ is a Nash Equilibrium. Suppose for Player $i$, the mixed strategy $\sigma_i^*$ assigns strictly positive probability to two pure strategies, $s_i^1$ and $s_i^2$ (i.e., they are in the \textbf{support} of $\sigma_i^*$).
  
  Then, the expected utility of playing either pure strategy against the opponents' equilibrium profile $\sigma_{-i}^*$ must be exactly the same:
  \[
    U_i(s_i^1, \sigma_{-i}^*) = U_i(s_i^2, \sigma_{-i}^*)
  \]
\end{thmbox}

\begin{notebox}
  \textbf{Corollary:} For any Nash Equilibrium $\sigma^*$, the expected utility $U_i(\sigma^*)$ is exactly equal to the expected utility of \textbf{any} pure strategy $s_i \in S_i$ that is played with strictly positive probability (i.e., $\sigma_i^*(s_i) > 0$).
  
  If one pure strategy yielded a strictly higher expected payoff than another, the player would just shift all their probability to the better one, contradicting that $\sigma^*$ is an equilibrium!
\end{notebox}

\end{document}